\chapter{Memory}\label{ch:memory}

A G17 program reaches memory through ordinary per-lane loads and stores. Each instruction names a
base (a device binding, threadgroup memory, the imageblock, or a register-held pointer), an index (a register or an
immediate), a byte displacement and an element width; the width of the access is fixed by the \emph{opcode}, with no width field.
The tensor unit's "tensor loads" are per-lane gathers of the same kind, at addresses that ordinary ALU instructions
compute, through the same load and store machinery and the same eight-slot scoreboard (MM~\mm{7};
\cref{sec:tensor-memory-path}).

Every memory access has an \emph{address
contract} (which bytes are touched) and an \emph{arrival contract} (when the value exists in a register or in memory),
and the hardware checks neither. A wrong binding rank stores nothing and returns the fill value; a load far past its
buffer returned 0 at every tested offset; a consumer that does not wait on a load's scoreboard slot reads the register's previous contents; a
released imageblock coordinate collapses every lane to element 0; an undeclared imageblock reads 0; a mis-encoded
texture coordinate reads texel (0, 0). All of these completed with status 0 and no fault. This chapter calls a rule
safe only where a control could have failed it, and such controls showed that several fields first read as address
fields are arrival (scoreboard) fields. The
arrival contract is \cref{ch:synchronization}'s subject (\cref{sec:scoreboard}); this chapter covers the address side and the storage kinds.

Terms used below:

\begin{itemize}
\item \emph{Rn}: 32-bit general register n, named here by role ("R32, the per-lane index register"). Register-file rules,
  including the destructive \emph{release} lifetime (modifier value 16), are \cref{ch:register-file}.
\item \emph{\opn{0}} and \emph{\opn{4}}: the decoder's tags for the uniform register file and for the message staging file that texture
  messages read. A \emph{binding expression} is how the decoder prints a base: \code{bin(op0, const(K), 8)} for a device base,
  \code{bin(op0, const(K), 2)} for a threadgroup base; K is the binding constant and the last number the base's scale.
\item \emph{Binding rank}: a buffer's position in the binding list the program image declares (\cref{sec:address-computation}).
\item \emph{Operand 1}: the control word after a memory instruction's first operand. Bits 20–23 name the slot a producer fills
  (as slot + 1); bits 24–31 are the instruction's own wait mask (\cref{sec:scoreboard}).
\item \emph{Element code}: the access-width field of tensor and vector loads: 1, 2, 4, 8 or 16 bytes (MM~\mm{7}, MM~\mm{25.37}).
\item \emph{Immediate-address form}: a load or store whose index and displacement are both immediates (\cref{sec:immediate-address-vector-forms}).
\item "Below Metal" means a pure IOGPU submission with no AGXMetal driver loaded
  (\code{docs/g17-first-dispatch-trials.md}; the interface is \cref{ch:submission-below-metal}).
\item Encoding and field maps are \cref{ch:instruction-encoding}; metadata and the constant-program format \cref{ch:kernel-metadata-image}; bandwidth, latency and
  instruction-fetch numbers \cref{ch:performance}.
\end{itemize}

\section{Address computation}\label{sec:address-computation}

A register-indexed access forms its address as follows (measured on the 128-bit vector load and through the emulator's
byte-for-byte agreement with shipped dispatches; the other source shapes below are what Apple's compiler emits):

\begin{lstlisting}
address = base + index_register x element_width + displacement        (width and displacement in bytes)
\end{lstlisting}

Element width scales only the index; the displacement is unscaled. A constant address has no index register, width 1 and a byte
displacement. The immediate-address family (\cref{sec:immediate-address-vector-forms}) uses a different formula, and tensor addresses are this
formula applied to an index register that ordinary arithmetic computes (below).

Two project files give conflicting answers on whether the displacement exists and what it counts.
\code{isa/g17-addressing-rules.json} compiled \code{C[tid] = A[tid + k]} with \code{uint tid} for six element types and found that no operand of the load moved with k: Apple's compiler instead added k to the index with \op{10279}, and k = 1, 2, 4, 8 encoded as immediate 1, 2, 4, 8 for char through long. Its conclusion was "the
load carries NO displacement field" and \code{address = base + (index + displacement) x sizeof}. \code{isa/g17-address-formula.txt}
compiled \code{f[i + 7]}, \code{w[i + 7]} and \code{L[i + 7]} and found a displacement operand carrying 28, 14 and 56: bytes, not
elements. Both observations are correct for the source shapes they compiled. On hardware, setting the 128-bit vector load's displacement to 4, 8, 12 and 20 moved the access by exactly that many bytes in 10 of
10 runs (MM~\mm{25.141.18}), and the repository's emulator, which models every register-indexed access as
\code{index x element + displacement} in bytes, agrees byte for byte with the GPU on 316 of 316 and 172 of 172 shipped
dispatches (MM~\mm{25.197}).

Compile-side and decoder facts about the same operands:

\begin{itemize}
\item A constant address: \code{u[1]} is 4, \code{u[7]} 28, \code{u[64]} 256 for uint (Apple's compiler emits,
  \code{isa/g17-address-formula.txt}).
\item For \code{A[tid + k]} with an unsigned 32-bit index, Apple's compiler puts k into the index with the add immediate (1 to
  255) or, at 256 and above, \op{10282} or \op{10283}, and leaves the
  displacement 0 (Apple's compiler emits, \code{isa/g17-addressing-rules.json}). The file does not record why; a reader may
  not assume the displacement can replace that add when an unsigned index sum may wrap.
\item The add immediate's 8-bit immediate is scattered across the word (value bits 0–1 at byte1{[}0:2{]}, 2–4 at byte3{[}5:8{]}, 5 at
  byte5{[}7{]}, 6–7 at byte8{[}0:2{]}); the map predicted 15 never-compiled values byte for byte (decoder, same file).
\end{itemize}

\emph{Correction:} "the load carries no displacement field" (\code{isa/g17-addressing-rules.json}) holds only for its unsigned
\code{tid + k} sweep. The field exists and counts bytes (MM~\mm{7}, MM~\mm{25.141.18}, \code{isa/g17-address-formula.txt}).

Rules a compiler has to follow (Apple's compiler emits, \code{isa/g17-address-formula.txt},
\code{isa/g17-addressing-rules.json}):

\begin{itemize}
\item \emph{No index-scale field.} \code{u[i * 4]} emits the same load as \code{u[i]}, with the multiply in a separate ALU instruction.
\item \emph{Width is the opcode.} char, short and half use \op{12646} and \op{17193}; int and float use
  \op{12682} and \op{17229}; long uses \op{12691} and the store
  \op{17238} (no glossary entry); int4 uses \op{12709} and \op{17256}.
\item \emph{Index width is the opcode too.} A ulong index selects \op{12688}, preceded by a
  64-bit shift pair that scales the index.
\item \emph{A runtime base plus index} (\code{A[j + tid]}, j loaded) is a load of j, a 32-bit add, then the indexed load.
\item \emph{Misalignment} does not select a form: an int at byte offset 1 through a char-pointer cast compiles to the same
  plain 32-bit load and store; whether hardware permits it for scalars is open. A vector load displaced by 4 bytes (not
  16-byte aligned) "lands whole" (measured, MM~\mm{25.141.18}).
\item No negative displacement was ever emitted (the compiler folds negative offsets elsewhere); the field's signedness is
  open.
\end{itemize}

The address is formed in 64 bits (measured, \code{isa/g17-predication.json}). A load of \code{A[1073741824 + tid]}, $2^{32}$
bytes past a 256-word buffer, was chosen so that a 32-bit byte wrap would land back inside the buffer; it returned 0, not the
wrapped word (\cref{sec:out-of-range-accesses-observed}).

A device base's binding constant is K = 4 × the binding rank, and the rank is the buffer's
position in the list the image declares, not the Metal slot number: a kernel binding only buffer C makes C rank 0
(measured, \code{agxforge/g17/cc.py}; MM~\mm{0.8}). The list order is internal bindings ascending, then user bindings ascending;
reconstructed from the declaration it reproduces Apple's list in 10,964 of 12,044 corpus kernels (91.0\%), never
differing by order alone (Apple's compiler emits, \code{agxforge/g17/resource.py}). A wrong rank makes a store return its fill
value at status 0 (measured, same file). Below Metal, declaring one unused buffer between A and C changed exactly one
code byte, the store's descriptor operand, from 0x04 to 0x08; swapping two 64-bit entries of the argument block
swapped which buffer a load read; and an entry pointed at an unviewed offset of an allocation, or at a different
registered allocation, was followed by loads and stores alike (measured, \code{docs/g17-first-dispatch-trials.md},
"Authored device atomic below Metal" and "Input descriptor address selection" onward). Arbitrary addresses and other
allocation classes are untested.

The 128-bit vector load takes its 14-byte form when it carries bit 37 or a
displacement above 255 (decoder, MM~\mm{6}). Tensor loads carry up to 255 bytes in the 12-byte form and 32,767 in the
16-byte form; past 32,767 the offset is folded into a per-tile base register (Apple's compiler emits, MM~\mm{25.114.3}).
This compiler does the same, visible as a kernel growing from 74 to 80 instructions, with no failure out to 1 MiB of
leading dimension (measured, MM~\mm{15}). An identical offset-bearing body was bit-exact at stream positions 1 to 12,
after 0 to 4,096 bytes of padding and at offsets 2 to 40,002 bytes straddling the boundary; within that range there is no
stream-position constraint (233 + 193 dispatches, MM~\mm{15}). Tensor B offsets are measured from 2 to 47,104
bytes; beyond is outside the measured domain (MM~\mm{25.114.3}).

Tensor addresses are ordinary arithmetic (measured, MM~\mm{7}). Apple's library computes one per-lane index register,
R32, the element offset of the lane's first element in a row-major tile, and its eight tensor loads read
\code{[R32*2 + disp]} at byte displacements \{0, 2048, 32, 4096, 6144\} (eight rows, sixteen columns), each moving an 8-row
by 16-column sub-tile. This arithmetic is what produces the fragment layout; the index formula and the fragments are
\cref{sec:fragment-layouts}. The row stride has no field either: a non-default stride changes the mode immediate from 241 to 2289
and moves into the index (decoder, MM~\mm{7}); a stride built as \code{K*2} regardless of element size was wrong, with no
error reported, on 1,024 of 1,024 elements of a 32×32×64 int8 product (measured, MM~\mm{17} rule 6; \cref{sec:tensor-memory-path}).

\section{Load and store forms by width}\label{sec:load-store-forms}

\Cref{tab:memory-forms} lists, by width, the forms this chapter uses; \cref{app:c} lists every memory opcode. Lengths are those observed in Apple's 6,594-program
corpus (\code{isa/g17-wait-mask-census.json}), with lengths the cited section adds. Operand orders are given where the forms
are explained: the plain 32-bit store's are value 0, operand 1 at 1, index 5 and index lifetime 6 (\cref{sec:scoreboard}), the
threadgroup forms' in \cref{sec:threadgroup-memory} and the imageblock forms' in \cref{sec:imageblock-per-core-tile}.

\begin{xltabular}{\linewidth}{@{}>{\hsize=0.88\hsize}Xll>{\hsize=1.18\hsize}XX@{}}
\caption{Load and store forms used in this chapter, by width, with observed lengths in bytes and standing.}\label{tab:memory-forms}\\
\toprule
Glossary name & Opcode & Lengths & Width and notes & Standing \\
\midrule
\endfirsthead
\tablecontinued{5}
\toprule
Glossary name & Opcode & Lengths & Width and notes & Standing \\
\midrule
\endhead
plain 32-bit load & \opn{12682} & 8, 10, 14 & 4 bytes, element code 4 & measured (MM~\mm{25.197}; below Metal) \\
16-bit load & \opn{12646} & 8, 10, 14 & 2 bytes (half, bf16, short, char); optional high-half destination & measured (MM~\mm{25.141.3}) \\
plain 64-bit load & \opn{12691} & 8, 14 & 8 bytes into a pair, element code 8 & measured (glossary) \\
three-component vector load & \opn{12700} & 8, 14 & 12 bytes into a 3-register tuple; the access steps 16 & measured (MM~\mm{25.90}) \\
128-bit vector load & \opn{12709} & 8, 10, 14 & 16 bytes into a 4-register tuple, element code 16; byte displacement & measured (MM~\mm{25.141.18}) \\
immediate-address one-word load & \opn{12688} & 8, 10, 14 & word at (index + displacement) × element; mask 2066 & measured (MM~\mm{25.177}) \\
tensor load, half/bf16 & \opn{12674} & 12, 16 & two words (four half/bf16), element code 2; immediate element mask & measured (MM~\mm{7}) \\
masked tensor load, half/bf16 & \opn{12675} & 12, 16 & as \opn{12674}, with a 16-bit mask register & measured (MM~\mm{7}) \\
int8 tensor load, one word & \opn{12656} & 12, 16 & one word (four int8), element code 1 & measured (MM~\mm{7}) \\
plain 32-bit store & \opn{17229} & 8, 10, 14 & 4 bytes (sub-form 00) & measured (MM~\mm{25.117}, MM~\mm{25.145.4}) \\
128-bit vector store & \opn{17256} & 8 (corpus) & 16 bytes from a 4-register tuple; uses every wait bit 24–31 & measured (MM~\mm{6}) \\
store, sub-form 01, texel/word-slot & \opn{17235} & 8, 10, 14 & one word at an immediate address; reads a texture fetch; mask 2066 & measured (MM~\mm{25.62}) \\
two-component store & \opn{17244} & 8, 10, 14 & words k and k+1 from a pair; mask 2098 & measured jointly with \op{17262} (MM
\mm{25.177}) \\
tensor store & \opn{17257} & 10, 12, 16 & four words; immediate mask & measured (MM~\mm{6}, MM~\mm{7}) \\
masked tensor store & \opn{17258} & 10, 12, 16 & four words; mask register, bit j enables word j & measured (MM~\mm{7}) \\
threadgroup-memory load & \opn{12364} & 8, 14 & 2 or 4 bytes & measured (\cref{sec:threadgroup-memory}) \\
threadgroup-memory store & \opn{13288} & 8, 10, 14 & 2 or 4 bytes & measured (\cref{sec:threadgroup-memory}) \\
imageblock load, 32-bit & \opn{12151} & 14 & one 32-bit member & measured (\cref{sec:imageblock-per-core-tile}) \\
imageblock store, 32-bit & \opn{13075} & 14 & one 32-bit member & measured (MM~\mm{25.96}) \\
\bottomrule
\end{xltabular}

Also used, outside \cref{tab:memory-forms}: \op{12673}, which Apple's compiler emits to load pointer +
immediate from a register-pair base; \op{12073}, a runtime 32-bit address for
thread-private arrays (Apple's compiler emits, MM~\mm{25.39.2}); \op{12692}, element code 4 (measured,
MM~\mm{25.37}), and the masked twins of the int8 and fp32 tensor loads; 16-bit store (\opn{17193}, lengths 10 and 14, measured,
glossary) and the 64-bit store \op{17238} (no glossary entry; lengths 8 and 10; Apple's compiler emits); \op{17202}, measured in the packed int8 epilogue (MM~\mm{25.105}); the rest of the immediate-address family (section
7.3); the spill stores (\cref{ch:register-file}). Two forms seen only in MLX's quantized matmul have an open address space: \op{12352} and \op{13276} (MM~\mm{25.144.1}). Decoder-only
and unmodelled forms are in MM~\mm{25.88} and MM~\mm{25.192}. The repository's encoders round-trip 1,145 of 1,190 and 3,419 of 3,793 loads
and stores; the unclaimed sub-forms are named in MM~\mm{15}. 132 memory opcodes declare implicit uses of \code{SR\_LOCAL\_X} and
\code{SR\_LOCAL\_Y} (\op{12118} to \op{12135} among them, no glossary entry for either); none occurs in the shipped corpus (decoder,
\code{isa/g17-local-indexed-memory.txt}).

\section{Immediate-address and vector forms}\label{sec:immediate-address-vector-forms}

The immediate-address family is measured in MM~\mm{25.177}. Its operands, in the order of \op{17235}, are value or destination, operand 1,
mask, binding, 0, index, displacement, element; the address is the word at \code{(index + displacement) x element} of the
binding, with index and displacement both immediates. The mask field is a component mask in bits 4–7 above 0x802:
2066 one word, 2098 two, 2162 three, 2290 four (bits 4–7 = 1, 3, 7, 15), filled from the tuple's lowest register up;
2065 is the halfword form (MM~\mm{25.195}). Members: \op{12688}, immediate-address two-word
load (\opn{12697}, MM~\mm{25.184}), \op{12706}, \op{12715}, \op{17244}, \op{17253}, \op{17262}, \op{17199}. \code{isa/g17-addressing-rules.json} proposed the same
\code{(index + displacement) x element} equation for the register-indexed form; it holds for this family and not for that one.

\emph{Worked example: a wrong-width model of this family} (emulated, then executed on hardware; MM~\mm{25.177}). The first model of
this family loaded each word with element size 1, keeping only its low byte, and still matched Apple's kept GPU hash on
9 of 10 census sources, because the census inputs are small integers, (7i + 3) mod 61. Only the one source with
float data failed; it read as zero (5.0 is 0x40a00000). Rerun with full-width random buffers (full-mantissa floats with
exponents to $2^{\pm20}$, uints over all 32 bits), all 10 sources match Apple's GPU byte for byte, and each deformed model is
refuted per form (sources agreeing, then sources where a low-byte load, a reversed load and a first-only store differ):
one-word load 5 / 5, 3, 3; four-word load 6 / 6, 6, 6; four-word store 8 / 6, 6, 8; two-component store 2 / 1, 1, 2,
receipted only jointly with the four-word store. The same protocol later receipted the three-word store (6 of 6,
MM~\mm{25.192}) and the 16-bit store at an immediate address (3 of 3, MM~\mm{25.195}). A memory form is marked
measured in this reference only if its test data could have exposed a wrong width.

In constant programs the base is a register pair holding a binding pointer that \op{14061} fetched; the address is pointer plus byte offset (measured, MM~\mm{25.179}; \cref{sec:constants-uniforms}).

The vector forms:

\begin{itemize}
\item \op{12709}: byte displacement (\cref{sec:address-computation}). Apple's compile of MLX's qmv carries 16, 32, 48 for
  \code{x[k+4]}, \code{x[k+8]}, \code{x[k+12]}. \emph{Correction:} \code{asm.encode\_vec4} calls the field "units of eight bytes"; for loads it is
  bytes (MM~\mm{25.141.18}). The store's units, witnessed only at 0, 16 and 32, are open.
\item Tensor loads and stores move popcount(mask) consecutive elements into the tuple, lowest address in the lowest half
  (measured, \code{tools/g17emu.py} EVIDENCE{[}12674{]}). int8 operands need \op{12656}, not the
  generic four-word load (measured, MM~\mm{7}).
\item This compiler, measured: vectors placed past a 32,767-byte displacement have their base folded into the index
  register once, and two per-channel vectors must lie within one field of each other (MM~\mm{25.165}); with loop-carried
  addresses (\code{ptr\_addr}), later loads read one carried index at displacements 16, 32, 48 and instructions per trip fall
  121 to 109, 184 to 171, 69 to 59 and 309 to 295 in four qmv loops, bit-exact against a 0x7f sentinel (MM~\mm{25.141.18}).
\end{itemize}

\section{Masked and predicated access}\label{sec:masks-predicated-memory}

\begin{itemize}
\item \op{612}: \code{bit j of dst = [lo <= src + j < lo + w]} for j in 0..3, unsigned. The second
  immediate is a \emph{width}. It was first read as an upper bound because nearly every witness had lo = 0, where the two
  readings coincide; two cases that separate the readings falsified the bound (measured, MM~\mm{7}). lo and w are 8-bit, so lowerings compute
  tile-relative predicates.
\item \op{621}: dst is a mask of \code{min(max(W - src, 0), 4)} low ones, W = operand 5, measured at
  W = 16, 32, 48 and 64 (src W-3..W gives 7, 3, 1, 0); the cap of four did not move with operand 3 at 8, 16, 32, 56
  (\code{isa/g17-address-formula.txt}). \emph{Correction:} the earlier fit \code{(0xFFFFFFFF >> src) \& 0xF} is the W = 32 case
  only. Every dispatch behind it had been authored on one witness whose operand 5 was 32, which coincides with the fitted
  saturation point. The source of the cap of four is open; there is no four-bit field.
\item \op{12675} and \op{17258}: mask bit j enables word j of the lane's
  four-word access; bits above 3 are ignored; unwritten words keep their contents; only the declared M × N region is
  written and poisoned padding never leaks (measured, six patterns across 32 lanes, MM~\mm{7}).
\item Ordinary predicated access (measured, \code{isa/g17-predication.json}): a masked-off load suppresses the \emph{value} (inactive
  lanes keep their sentinel); a masked-off store suppresses the \emph{access} (seven masks); effects are per lane and nested
  masks intersect. Whether a masked-off load still issues its access is unknown: the intended observable was a fault,
  and no out-of-range load faulted (\cref{sec:out-of-range-accesses-observed}). Apple's compiler emits predicated
  memory branchlessly between exec-mask instructions, with a skip branch for two-armed effects (\cref{ch:control-flow-execution}).
\end{itemize}

\section{Out-of-range accesses}\label{sec:out-of-range-accesses-observed}

The masking experiment needed a load that faults, to separate a suppressed access from a suppressed value (measured,
\code{isa/g17-predication.json} \code{wild\_address}). Unconditional \op{12682} reads at
4,194,304, 67,108,864, 268,435,456 and 1,073,741,824 words past a 256-word buffer (16 MiB to 4 GiB) ran through Metal,
each in a fresh process so that a device loss would cost only that process. All completed with status 0 and returned 0.
An in-bounds control returned its data; at the $2^{32}$-byte offset the load returned 0, not the in-buffer word a 32-bit
wrap would select. No fault occurred, so the masking question is still open.

Out-of-range store behaviour has not been measured. A timing driver once read and wrote up to 2.78~MB
outside its buffers, with uncharacterised effect (MM~\mm{17} rule 34). The emulator refuses a store outside a buffer, a model
bound (\code{tools/g17emu.py}).

Observed faults:

\begin{itemize}
\item Through Metal, one BIF0 read page fault came from a \code{regprobe\_multi} tensor-kernel probe whose configuration was not
  recorded; a different process, sized from a sibling kernel's layout, hung the GPU for about 320 s and the machine
  rebooted (MM~\mm{17} rule 34, as corrected; \cref{app:a}).
\item Below Metal, malformed launch state produced BIF0 read page faults at unmapped addresses while Submit still returned
  0, both callbacks ran and the recovery counter advanced: a guessed 2,048-byte threadgroup resource word (fault at
  address 0; the matched word then ran), an all-zero allocation 0 (fault at 0x10000010000, its end), an oversized
  allocation layout (0x100000d0000, 0x100000d8000) and a malformed code-placement packet (0x80) (measured,
  \code{docs/g17-first-dispatch-trials.md}). The requester of those reads was not identified.
\end{itemize}

\emph{Correction:} MM~\mm{17} rule 34's "only the accelerator's own reads have faulted" is scoped to work submitted through
Metal; below Metal, launch-state errors fault too.

Two silent failures resemble range errors. Undersizing a host dispatch's completion scope (\code{dim\textasciicircum{}2 * esz})
returns wrong or zero data past that offset with no error (measured, MM~\mm{17} rule 5). A dead load landing late can
overwrite a reused address register so that the next load reads out of range and returns 0 (\cref{sec:dependences-issue-scalar}, MM~\mm{25.114}).

Through Metal, a stray plain 32-bit load inside the tested range returns 0 without
faulting. A reader may not assume: that other load forms, resource kinds or offsets behave the same; that an
out-of-range store is dropped; that reading an unmapped address is harmless below Metal; or that a 0 means "out of
range", since an unwaited load commonly reads 0 too (\cref{sec:scoreboard}). Keeping accesses in bounds is left to the program. Threadgroup accesses
past 32,768 bytes and imageblock coordinates outside the tile are refused by the emulator only as a model bound; one
imageblock read at x = 32, outside a 32-wide tile, completed and was not scored (MM~\mm{25.104.2}).

\section{Threadgroup memory}\label{sec:threadgroup-memory}

Size limits:

\begin{itemize}
\item 32,768 bytes per threadgroup is Apple's compiler's hard refusal; 48, 60 and 64 KiB are refused and 96 KiB drops the
  compiler connection (measured, MM~\mm{7}). Metal itself does not check it for an authored container: a container
  declaring 1 MiB of static threadgroup memory built a pipeline that reports \code{staticThreadgroupMemoryLength = 1048576}
  (measured, MM~\mm{25.147}, nothing dispatched). What the hardware does with an oversize declaration is unmeasured.
\item 57,344 bytes (56 KiB) per core, inferred from a residency step: with 128 threads per threadgroup and 4,096
  threadgroups, a static request of 28,672 bytes ran in 0.00304 s and 28,688 bytes in 0.00382 s, so two threadgroups
  fit at exactly half the budget and one above it (measured step, MM~\mm{25.49}). Residency rules are \cref{ch:device-execution-model}.
\item \code{maxTotalThreadsPerThreadgroup} stays 1,024 across threadgroup memory from 64 B to 32,768 B (measured, MM~\mm{25.47});
  \code{staticThreadgroupMemoryLength} is the value the kernel's metadata declares (MM~\mm{25.147}).
\end{itemize}

The forms are \op{13288} and \op{12364}, plus a constant-index threadgroup
load \op{12367} (no glossary entry) (Apple's compiler emits, \code{isa/g17-threadgroup-pair.txt}). Their operands are value or
destination, operand 1, width, base, 0, index, index lifetime, displacement, element. The base's binding expression
has scale 2 where device forms use 8 (notation list), with 0 exceptions in 25,553 corpus memory instructions
(decoder, \code{isa/g17-address-formula.txt}). The address is base + index x element + displacement (the emulator's model,
measured on its receipts, \code{tools/g17emu.py} EVIDENCE{[}13288{]}).

\emph{Worked example (compiled by Apple, decoded; the same shape executed below Metal).} From one kernel, so the two
addresses agree by construction (\code{isa/g17-threadgroup-pair.txt}, \code{a6-tgidx}):

\begin{lstlisting}
threadgroup-memory store (op13288), 10 bytes  0f040302030611c00020
    operands: R105 (the stored value), operand 1 = 0x80000010, 2066, bin(op0,const(4),2), 0,
              index register, 16, 0, 4      (op0: uniform register file)
threadgroup-memory load (op12364), 14 bytes   0f0403021a2610c0410080000000
\end{lstlisting}

The store's operand 1 is 0x80000010: bit 31 set, a wait on slot 7, the slot this kernel's loads fill (\cref{sec:scoreboard}); the
trailing 4 is the element width. Below Metal, a program of this shape, each lane storing \code{A[i]} to scratch, a
threadgroup barrier, then lane i reading slot \code{i XOR 32}, produced 64 of 64 values \code{11 * (i XOR 32) + 5} across two
simdgroups, and a 2,048-byte allocation held 512 separately addressable words with every written value recovered
(measured, \code{docs/g17-first-dispatch-trials.md}, "Parallel static-threadgroup launch mapping"). The static size is a
launch resource word, not an instruction field (\cref{ch:submission-below-metal}): a guessed value for 2,048 bytes faulted, and the matched
value Metal writes (\code{0f00080c00000000}) does not follow the simple size/8 pattern of the smaller sizes.

Further measurements:

\begin{itemize}
\item The threadgroup region is provisioned by an ordinary threadgroup load or store: removing the one threadgroup store
  made a threadgroup atomic write nothing (measured, ledger \code{g17-the-threadgroup-atomic-subtracts-and-needs-a-region}).
\item The threadgroup load's bit sensitivity: clearing byte4{[}3{]}, byte4{[}4{]} or byte6{[}4{]} made the read return 0;
  byte5{[}5{]}, byte8{[}6{]} and byte10{[}7{]} are inert; of the three-bit variants, 000 read correctly and 010, 100 and 110 read 0
  (measured, \code{isa/g17-execution-tgbits-results.json}, \code{-tgtrio-results.json}). These are the positions of the 128-bit
  vector load's fill-slot field (MM~\mm{6}); that they are a slot field here fits, but is not established.
\item A cross-simdgroup combine (word \code{sg} per simdgroup, one threadgroup barrier, an ascending fold) is exact for S = 2 and
  4, sum and max (measured, MM~\mm{25.90}). Ordering between simdgroups needs the barrier (\cref{sec:barriers-fences}); hand-off costs are
  \cref{ch:performance}.
\end{itemize}

\section{Imageblock (per-core tile memory)}\label{sec:imageblock-per-core-tile}

The imageblock is a per-threadgroup tile indexed by thread position and shared within the threadgroup (measured,
\code{isa/g17-imageblock-execution.txt}): 32 lanes wrote 32 distinct values to their own pixels, and every lane that read pixel
(0, 0) wrote lane 0's value into a buffer prefilled with 0xdeadbeef.

The forms are \op{13075} and \op{12151}, with operands
\code{[value, operand 1, width 18, member byte offset, 1, coordinate register, coordinate lifetime, dx, dy]} (measured,
MM~\mm{25.145.2}). The coordinate register packs x in its low 16 bits and y in its high 16; Apple's compiler builds it with
\op{10286}. The declared pipeline allocated 128 bytes for 32 lanes of one 32-bit member
(MM~\mm{25.104.1}).

\textbf{Worked example: the imageblock store's wait bit} (eleven rounds; compiled, decoded and executed on hardware;
MM~\mm{25.96}, MM~\mm{25.104}). The target was a tensor GEMM, every byte from this compiler, that stages its output through
the imageblock. Three defects gave the same symptom, a neighbour read returning the reader's own value or 0.

\begin{enumerate}
\def\labelenumi{\arabic{enumi}.}
\item \emph{The released coordinate.} The store released its coordinate register (operand 6 = 16) while the following read
  reused it. Release is destructive (\cref{ch:register-file}): a one-bit same-run control, keep (0) against release (16), collapsed
  all 32 lanes to element 0 (measured, MM~\mm{6}; \code{isa/g17-imageblock-dseries-receipt.json}). Fix: operand 6 from liveness.
\item \emph{The dx operand.} With the coordinate kept, the x offset (operand 7 = 1) still had no visible effect in a tensor
  program. Apple's compiler never uses it for a neighbour: it computes \code{(tid.x + 1) \& 31} with \op{10289} and \op{435} and reads with operand 7 = 0. This compiler now addresses the
  neighbour by an explicit column (measured, MM~\mm{25.104.2}).
\item \emph{Operand 1.} Apple's own tensor kernel then read its neighbour on all 32 lanes and this compiler's did not. The only
  field-level difference was the store's operand 1: Apple's 0x80000010 against this compiler's 0x01000010. Swapping it
  alone moved the result both ways, and bit 31 decided: 0x80000010 and 0x81000010 passed, 0x00000010 and 0x01000010
  failed. The reading recorded at the time was that bit 31 selects storage other lanes can read.
\end{enumerate}

A later experiment refuted that reading. Bits 24–31 of the imageblock store's operand 1 are its \emph{wait mask}
(bit 24 + s waits on slot s), the same field general stores and MMAs carry (\cref{sec:scoreboard}). In a census of all 57 corpus
imageblock-store instances the bits equal exactly the slots the store's sources are pending on: 47 need none and carry none, 8 carry bit
24 for the coordinate's special-register reads, 2 carry bit 25. A prediction made before compiling, that an Apple kernel
storing a freshly loaded value (its load fills slot 7) would carry bit 31, was confirmed. \Cref{tab:imageblock-wait-arms} gives
the hardware runs of that kernel (\code{h[t] = 5t + 3}, a device barrier, \code{v = h[t \textasciicircum{} 1]}, an own-pixel store of v, read back), three per arm.

\begin{xltabular}{\linewidth}{@{}>{\hsize=1.19\hsize}X>{\hsize=0.81\hsize}X@{}}
\caption{Imageblock store operand 1 against lanes that read their neighbour's value correctly, in Apple's kernel that stores a
freshly loaded value. Three runs per arm.}\label{tab:imageblock-wait-arms}\\
\toprule
store operand 1 & lanes correct \\
\midrule
\endfirsthead
\tablecontinued{2}
\toprule
store operand 1 & lanes correct \\
\midrule
\endhead
0x80000010, as compiled (waits slot 7) & 32 / 32 \\*
0x10, wait cleared & 0 / 32, every lane 0 \\*
0x2000010, waits slot 1, which nothing fills & 0 / 32, every lane 0 \\*
\bottomrule
\end{xltabular}

Without the right slot the store takes its value register before the load writes it; the failure is in the store,
not in the neighbour read. Earlier deformations of operand 1 in Apple kernels had found the bits inert because those kernels stored
ALU results, which need no wait. Round 11 routed the stored value through a waiting copy and kept the slot-0 wait for
the coordinate (0x01000010); every arm was bit-exact (MM~\mm{25.104.4}). A store's bits 24–31 must name every slot its
value and coordinate registers are pending on.

\emph{Correction:} the "bit 31 selects coordinate-addressed staging" wording in MM~\mm{25.89}'s round-9 entry, in MM~\mm{25.145.2}
and in the \code{tools/g17emu.py} comments on the imageblock store and load is the superseded reading; bit 31 is the wait on slot 7 (MM~\mm{25.96},
MM~\mm{25.104.3}). There is no per-thread versus shared distinction in the evidence. The emulator's admission of cross-lane
reads only for words stored with bit 31 is a conservative restriction of the model; no hardware result calls for it.

Declaration and API requirements (measured): an undeclared image reads 0 even when the tile is sized at encode time with
\code{setImageblockWidth}: C{[}0, 0..31{]} were all 1.0, the read having returned 0, in three queries (MM~\mm{25.104}). The pipeline
must come from \code{newComputePipelineStateWithDescriptor:}; a function-name pipeline reports
\code{imageblockMemoryLengthForDimensions:} as 768 bytes at 32×1; \code{setImageblockWidth:height:} must match the threadgroup
(\code{isa/g17-imageblock-execution.txt}). The declaring metadata is \cref{ch:kernel-metadata-image}'s.

The imageblock barrier is \op{447} at scope 532, bytes \code{475100060000}
(\code{threadgroup\_barrier(mem\_threadgroup\_imageblock)}), which Apple's own-element kernel carries; Apple's compiler elided a
source \code{simdgroup\_barrier(mem\_threadgroup)} between a one-simdgroup imageblock store and load (Apple's compiler emits,
\code{isa/g17-barrier-scopes.txt}, \code{isa/g17-imageblock-execution.txt}).

\textbf{Not known.} Use by more than one simdgroup or threadgroup (never witnessed, MM~\mm{25.145.2}); the bit-4 flag (0x10)
every store's operand 1 carries (MM~\mm{25.96}); the read's operand-1 word Apple emits, 0x2000800000 (MM~\mm{25.104.2}). Staging a
fragment through the imageblock is bit-exact but frees no registers; its cost is in \cref{sec:chaining-feed-modes} (MM~\mm{25.106}).

\section{Constants and uniforms}\label{sec:constants-uniforms}

A constant is an immediate inside the instruction, an entry in the \emph{constant pool}
(compile-time data in metadata slot 13 of the container; Apple's compiler emits, \code{agxforge/g17/cooperativemetadata.py}), or a
\emph{published uniform}, a value Apple's compiler computes per dispatch in a separate \emph{constant program} that runs once before the main
program and writes the uniform register file. Main reads pool entries and published uniforms through the same operand
syntax, \code{bin(op0, const(K), W)}, so the operand alone does not say which one it is. Treating a published uniform as a
pool entry gave 6 silent wrong answers (measured, MM~\mm{25.179}).

The constant program was measured in MM~\mm{25.179}:

\begin{itemize}
\item The emulator reproduces Apple's GPU output by running the constant program first, as one thread over the same buffers,
  and handing its uniform writes to main: 6 of 6 sources byte-identical on full-width random buffers; three deformed
  models (every published uniform reads 0; the pointer names the next binding; a 32-bit uniform lands with halfwords
  swapped) each differ on all 6.
\item Binding pointers: in all 21 constant programs studied, \op{14061} at halfwords 8 + K
  and 10 + K reads binding K's 64-bit pointer, low word then high. With K = 4 × rank (\cref{sec:address-computation}) this is the same rule
  as MM~\mm{25.59}'s "slots 8 + 4k and 10 + 4k for the k-th buffer". The pointer is then the register-pair base for the
  immediate-address loads (\cref{sec:immediate-address-vector-forms}).
\item Uniform destinations are main-program operations with one extra mode operand: \op{11185} is \op{11179}; \op{592} is \op{586}; \op{1038} is \op{1006}.
  A uniform destination must be written with one value across the active lanes.
\item Of 15,256 corpus constant-program uniform writes, 14,814 are read by main; the 442 never-read writes are all uniform /
  staging writes and unexplained (decoder, MM~\mm{25.120}).
\end{itemize}

In main, \op{428} takes its mask from a pool operand (\cref{sec:integer-arithmetic}), and the
shipped decode replay refused until the pool was loaded, at 768 uses (measured, MM~\mm{25.197}). \op{435} takes an 8-bit immediate; a wider mask such as 0xFF0 goes to the
uniform-operand \op{437} (Apple's compiler emits, MM~\mm{25.192}).

In the uniform / staging write (decoder, MM~\mm{25.120.1}), the value the decoder printed as a
"publish index" is a scoreboard slot + 1. With a uniform destination it is 0 (47,269 of 47,269); with a destination in the message staging file
(\opn{4}) it is 1–8 (131,313 of 131,313). The reader of a staging write is a message instruction (\op{14661}, \op{10909}, \op{15813}), and it waits on slot index - 1: of 189,898 staging writes,
174,909 are waited on by the reader and 14,989 by an intervening instruction; none by neither. A staging write is a
late value, like a load (\cref{sec:scoreboard}).

The binary format of the constant program and pool is \cref{ch:kernel-metadata-image}'s.

\section{Textures}\label{sec:textures}

The texture forms are \op{15813}, \op{14661}, \op{10909} and \op{11564} (glossary). No operand carries the coordinate. Two reads at different coordinates are byte-identical, and the
coordinate is published into the message staging file (\opn{4}) at \code{[op4+0*4]} and \code{[op4+2*4]} before the message (measured,
\code{tools/g17texrun.py}).

An authored read, \code{A[400] = TX.read(x, y).x}, was executed (\code{tools/g17texrun.py},
\code{isa/g17-execution-texture-results.json}) on a texture holding y*1000 + x + 7, chosen so a wrong texel, a wrong row, a broadcast coordinate and an unbound texture
give four different wrong answers: coordinate (5, 1) gave 1012, (2, 3) 3009, and lane 31 alone at (7, 3) 3014; all 32
lanes racing for one word gave an ambiguous 3007. The same runs showed that the texture operand is a dense index over
the textures the function uses, not the Metal binding index (five constructed kernels); and publishing a register
coordinate with a uniform / staging write needs operand 2 = 1048576; with 16777216 every lane silently reads texel (0, 0).

Which store reads a texture fetch is unresolved. \op{17235} reads a texture fetch (measured,
MM~\mm{25.62}), and two-component store (\opn{17244}, 14 bytes) read one once (texel 201 at (1, 2), one 4×4 R32Uint texture, one
thread; MM~\mm{25.146}). For \op{17229}:

\begin{itemize}
\item A dispatched bisection of Apple's reading program found that swapping only its store for this compiler's 8-byte
  plain store stopped the read; that store, and a 12-byte ALU copy placed after the fetch, saw the destination register's
  previous contents, at zero and at six intervening instructions. The ledger concluded "not timing" (ledger
  \code{g17-only-one-consumer-can-read-a-texture-fetch}).
\item Apple's own \code{tex2d-read} stores the texture read's result with an 8-byte plain store whose operand 1 is 16777232
  (0x1000010, a wait on slot 0), and that read publishes on slot 0 (Apple's compiler emits, MM~\mm{25.146},
  \code{tools/g17applestorefetch.py}). This compiler's encoder writes no wait on the 8-byte plain store, and the ALU copy in the
  bisection waited on slot 7, not slot 0.
\item The bisection's arms are consistent with a missing wait on the fetch's slot; "not timing" does not
  rule out a wait, since unwaited late values stayed wrong after 64 intervening instructions elsewhere (\cref{sec:scoreboard}). The settling
  dispatch, Apple's \code{tex2d-read} run and read back, has not been made.
\end{itemize}

A fetch result can safely be treated as a late value whose consumer waits on the fetch's slot, and the sub-form 01
store can be relied on to read it. Do not assume that the plain 32-bit store cannot read a fetch, or that distance
substitutes for the wait.

Pixel format is descriptor-side (Apple's compiler emits, \code{isa/g17-texture-sampler-surface.json}):
\code{texture2d<uint>} and \code{texture2d<float>} (R32Uint, R32Float) compile to identical code and metadata; R16Float changes
only the half-flavoured instructions. A sampler adds one register and a 4-byte sampler-state descriptor and swaps the
texture read for the texture sample. Texture metadata slots are \cref{ch:kernel-metadata-image}'s. Sampled and write texture
classes have no dispatched receipt (open).

\textbf{Safe assumptions.} The register-indexed law with unscaled byte displacements for the forms of
\cref{sec:address-computation}; the immediate-address law for its eight measured members; width and index width as
opcode choices; materialised strides; 64-bit address formation; and that every late value (load, imageblock read,
staging write, texture fetch) needs its consumer's wait.

\textbf{Not known.} Out-of-range stores; negative displacements; misaligned scalar access; the vector store's
displacement units; whether a masked-off load issues its access; what the threadgroup load's sensitive bits are;
multi-simdgroup imageblock use; the imageblock store's bit-4 flag; the address space of the two MLX-only forms
(\cref{sec:load-store-forms}); the texture-fetch consumer; sampled and write textures; the origin of range mask, 16-bit
source's cap of four; hardware behaviour for an oversize threadgroup declaration.

\textbf{Evidence.} MM~\mm{0.8}, MM~\mm{6}, MM~\mm{7}, MM~\mm{15}, MM~\mm{16}, MM~\mm{17} rules 5, 6 and 34, MM~\mm{25.47}, MM~\mm{25.49}, MM~\mm{25.62}, MM~\mm{25.89}, MM~\mm{25.90},
MM~\mm{25.96}, MM~\mm{25.104} (25.104.1 to 25.104.4), MM~\mm{25.105}, MM~\mm{25.106}, MM~\mm{25.114}, MM~\mm{25.114.3}, MM~\mm{25.120}, MM~\mm{25.120.1},
MM~\mm{25.141.18}, MM~\mm{25.144.1}, MM~\mm{25.145.2}, MM~\mm{25.146}, MM~\mm{25.147}, MM~\mm{25.165}, MM~\mm{25.177}, MM~\mm{25.179}, MM~\mm{25.184}, MM~\mm{25.192},
MM~\mm{25.195}, MM~\mm{25.197}; \code{isa/g17-address-formula.txt}, \code{isa/g17-addressing-rules.json}, \code{isa/g17-predication.json},
\code{isa/g17-threadgroup-pair.txt}, \code{isa/g17-imageblock-execution.txt}, \code{isa/g17-barrier-scopes.txt},
\code{isa/g17-wait-mask-census.json}, \code{isa/g17-texture-sampler-surface.json}, \code{docs/g17-first-dispatch-trials.md},
\code{tools/g17emu.py}, \code{tools/g17texrun.py}, \code{tools/g17applestorefetch.py}.
